\scanpage{99}
\margin{-4-}

\uline{of regular classes of} $G$.

In 7.4 an example of $N$ is given: in case $G$ is of type $\mathrm{SL}(r+1)$ we obtain one of the classical normal forms under conjugacy. This special case suggests the problem of extending the normal form $N$ from regular elements to arbitrary elements. In 7.1 it is shown that $N$ is a closed irreducible subset of $G$, isomorphic as a variety to affine $r$-space $V$, and in 7.9 (this is the main lemma concerning $N$) that if $G$ is simply connected and $\chi_i$ ($1\leq i\leq r$) denote the fundamental characters of $G$ then the map $x\to(\chi_1(x),\chi_2(x),\ldots,\chi_r(x))$ induces an isomorphism of $N$ on $V$. Then in \S8 the proof of 1.4 is given and simultaneously the following important criterion for regularity is obtained.

1.5 Theorem. \uline{If} $G$ \uline{is simply connected, the element} $x$ \uline{is regular if and only if the differentials} $d\chi_i$ \uline{are independent at} $x$.

At this point some words about recent work of B. Kostant are in order. In [3] and [4] he has proved, among other thing, the analogues of our above discussed results that are obtained by replacing the semisimple group $G$ by a semisimple Lie algebra $L$ over the complex field (any algebraically closed field of characteristic $0$ will serve as well) and the characters $\chi_i$ of $G$ by the basic polynomial invariants $u_i$ of $L$. The $\chi_i$ turn out to be considerably more tractable than the $u_i$. Thus the proofs for $G$ with no restriction on the characteristic are simpler than those for $L$ in characteristic $0$. Assuming both $G$ and $L$ are in characteristic $0$, substantial parts of 1.1, 1.2, and 1.3 can be derived from their analogues

\mspage{98}
